% LSPV-5 REPORT::PRIMECRT1::2026-10-02T08:45:00-07:00
\documentclass[11pt]{article}
\usepackage[margin=.88in]{geometry}
\usepackage{amsmath,amssymb}
\begin{document}
\section*{PRIMECRT1: prime/CRT Boolean semantics and provenance closure}
\textbf{STATUS: PROVED/EXACT local algebraic and closure results + COMPUTATIONAL controls. P versus NP remains OPEN.}

\section*{CORPUS REVIEWED}
The current repository tree (264 paths at survey time) was enumerated. This run reviewed root \texttt{main.tex}; AGENTS/CODEX/pnp-research governance; TASK-001 through TASK-007; verified MAIN 0816 and relevant historical synthesis; PRIME-LQP-1; BITRANK1; MATRIXCERT1; EXCESS1/EXCESSDUAL1; RAMCOMP1; TRACECLOSE1; current GROK/CLAUDE artifacts; HEC/VSC-HEC/G-PPC-HEC/CRS-PPC/CCSPC/RCCN/Search--G materials; governance notes; and active experiment/result paths. Repository history/search was used to locate prior provenance and reacquisition formulations.

The local container could not clone GitHub because network/DNS access was unavailable, so the preflight script itself could not run locally. Its source was read and its repository invariants are checked through the GitHub connector before completion.

\section*{DUPLICATION / CONTRADICTION AUDIT}
PRIME-LQP-1 already proved that ordinary Boolean contexts cannot distinguish semantically equivalent histories and that deterministic polynomial trace annotations are reacquirable by TRACECLOSE1/REACQ1. This report adds an exact CRT realization of successive Boolean logic, a provenance-semiring collision/negation audit, a polynomial size bound for positional prime codes, and a stronger search-closure theorem under \(P=NP\).

\section*{Theorem 1: exact prime/CRT Boolean algebra}
Let \(Q=\prod_{j=1}^{k}p_j\) be square-free and
\[
E_Q=\{e\in\mathbb Z/Q\mathbb Z:e^2=e\}.
\]
CRT gives \(\mathbb Z/Q\mathbb Z\cong\prod_j\mathbb F_{p_j}\), and each field has only idempotents \(0,1\). Hence
\[
E_Q\cong\{0,1\}^{k}.
\]
For idempotents define
\[
e\wedge f=ef,\qquad \neg e=1-e,\qquad e\vee f=e+f-ef.
\]
These are exactly Boolean AND, NOT, and OR coordinatewise.

For \(n\) variables choose \(k=2^n\) coordinates indexed by assignments. Every Boolean function \(F\) has a unique idempotent residue \(e_F\) satisfying \(e_F\equiv F(a_j)\pmod{p_j}\). Thus the proposed prime/mod-\(q\) multisign idea has an exact realization. It is, however, an algebraic encoding of ordinary semantics rather than provenance.

\section*{COMPUTATIONAL n=5 test}
Committed artifacts:
\[
\texttt{experiments/prime\_lqp/prime\_crt\_test.py},
\qquad
\texttt{results/raw/prime\_lqp/2026-10-02\_LSPV5\_PRIMECRT1.json}.
\]
Seed: 20261002. Exhaustive \(k=5\) testing verified all 32 idempotent states. The \(n=5\) truth-table representation used 32 prime coordinates and a 169-bit modulus. Five hundred deterministic random 32-bit state pairs passed exact NOT/AND/OR checks.

Double-negation, absorption, and De Morgan pairs had equal truth tables and equal CRT semantic residues but different positional-prime codes. NOT, AND-with-\(x_2\), OR-with-\(x_2\), and XOR-with-\(x_2\) continuations preserved semantic equality for every pair. This is COMPUTATIONAL confirmation of the exact theorem, not a lower bound.

\section*{Theorem 2: positional prime provenance remains compact}
For a history \(h=(e_1,\ldots,e_T)\), let
\[
G(h)=\prod_{i=1}^{T}p_i^{c(e_i)},\qquad 1\le c(e_i)\le C.
\]
Then
\[
\log_2G(h)\le CT\log_2 p_T=O(CT\log T).
\]
Thus polynomial \(T,C\) give a polynomial-length exact ordered prime code. The experiment reached about 90,166 bits at \(T=1000,C=8\), consistent with this bound. The huge numerical value is not a complexity resource, and factorization hardness is inadmissible because the generating computation already knows the factors.

\section*{Theorem 3: plain positive primes are not a faithful full provenance semiring}
Unique prime factorization faithfully labels a single multiplicative monomial. It does not faithfully encode semiring addition of alternative derivations. Under
\[
x\mapsto2,\ y\mapsto3,\ z\mapsto5,\ w\mapsto11,
\]
the distinct provenance expressions
\[
xy+z\quad\hbox{and}\quad w
\]
both evaluate to 11.

Negation is a sharper obstruction. Gr\"adel--Tannen provenance with negation uses dual indeterminates with \(p\bar p=0\). Positive ordinary primes cannot preserve this annihilation law. Therefore plain prime arithmetic is \textbf{FALSIFIED/RETIRED as an exact full provenance-semiring semantics}. CRT idempotents repair Boolean semantics, not derivational history.

\section*{Theorem 4: PRIME-FNP-CLOSE1}
Assume \(P=NP\). Let \(R(x,w)\) be polynomial-time decidable and polynomially balanced, and let \(g(x,w)\) be polynomial-time computable with polynomial output length. Put
\[
L_R=\{x:\exists w\,R(x,w)\}.
\]
Then there exists \(H\in FP\) such that for every \(x\in L_R\),
\[
H(x)=g(x,w_x)
\]
for some valid witness \(w_x\).

\paragraph{Proof.}
For a prefix \(u\), define
\[
E(x,u)\iff\exists v\;R(x,uv)
\]
with total witness length bounded by the polynomial balance. This predicate is in NP and hence in P under \(P=NP\). Query \(E(x,u0)\) and choose the next bit accordingly. Polynomially many prefix queries recover a valid witness, after which \(g\) is computed in polynomial time. \(\square\)

\paragraph{SAME\_REALIZATION consequence.}
If a downstream prime condition accepts any polynomially balanced NP-verifiable legal realization, then under \(P=NP\) a legal realization and its prime annotation are polynomial-time reacquirable. If it instead requires the exact original history of a deterministic polynomial SAT decider, TRACECLOSE1 applies by rerunning the decider. Generic prime binding therefore does not escape the current acquisition closure.

A surviving source-specific relation would have to (i) not factor through semantics, (ii) reject arbitrary NP-verifiable replacement realizations, (iii) not be a polynomial-time functional of a deterministic polynomial trace, and (iv) nevertheless be forced on every relevant SAT algorithm by a model-preserving theorem. No such theorem is currently present.

\section*{Literature and barrier audit}
Green--Karvounarakis--Tannen establish semiring provenance. Gr\"adel--Tannen extend provenance to negation with quotient semirings and dual indeterminates. Ahrens--Kassing--Giesl--Katoen give semiring weights to rewrite steps and derive quantitative rewriting analyses. These validate provenance as a formal derivation layer but do not transfer to arbitrary SAT lower bounds.

PRIMECRT1 explicitly algebrizes and relativizes. DEAD/COPY/DELAY change syntax provenance without semantics; INTERP/RECOMP regenerate ordinary traces. Horn, dual-Horn, 2-SAT, affine/Tseitin, bounded-treewidth, padding, renaming, and independent products are equally representable by CRT, so the representation cannot distinguish hard from easy families.

\section*{CORPUS RECONCILIATION}
\textbf{What changed relative to root MAIN.}
PRIME-LQP is narrowed. Prime labels remain useful as an exact SAME\_REALIZATION namespace for experiments. Prime magnitude/factorization, operator-prime products, deterministic trace labels, and polynomially balanced NP-verifiable replacement-witness labels do not provide the final hinge. The structural team target \(\operatorname{rank}_{\mathbb F_2}\Phi=\omega(m)\) remains higher priority and unchanged.

\textbf{Adopted:} positional prime labels as instrumentation; semantic/provenance quotient framing; TRACECLOSE1/REACQ1; dual-indeterminate provenance as an exact literature analogue.

\textbf{Rejected transfer:} database/rewrite provenance theorems do not force an arbitrary SAT solver to retain annotations.

\textbf{Falsified/retired in generic form:} prime magnitude/factorization as cost; commutative operator-prime products as ordered history; plain positive-prime arithmetic as full provenance semiring; deterministic polynomial trace prime labels as final hinge; polynomially balanced NP-verifiable replacement-witness prime labels as a contradiction to \(P=NP\).

\section*{Single next target}
Use event-specific primes only as instrumentation in one source-matched SU/vMS/RMV LQP experiment. Construct two histories with the same current semantic state and different prime histories. For every later legal distinguishing restriction \(D\), classify it immediately as:
\[
D\text{ factors through semantics},\quad
D\in\text{polynomially balanced NP search},\quad
D\text{ is deterministic-trace replayable},
\]
or genuinely outside all three. Only the fourth class survives as a provenance candidate. If the source-matched \(n=5\) search finds no fourth-class restriction, downgrade PRIME-LQP to instrumentation only.

\begin{thebibliography}{9}
\bibitem{GKT07} T. J. Green, G. Karvounarakis, V. Tannen, ``Provenance Semirings,'' PODS 2007, 31--40.
\bibitem{GT25} E. Gr\"adel, V. Tannen, ``Provenance Analysis and Semiring Semantics for First-Order Logic,'' 2025, arXiv:2412.07986.
\bibitem{AKGK25} E. Ahrens, J.-C. Kassing, J. Giesl, J.-P. Katoen, ``Weighted Rewriting: Semiring Semantics for Abstract Reduction Systems,'' LIPIcs FSCD 2025, 337:6:1--6:21.
\bibitem{HNO96} E. Hemaspaandra, A. Naik, M. Ogihara, A. Selman, ``P-Selective Sets and Reducing Search to Decision vs Self-Reducibility,'' JCSS 53(2), 1996, 194--209.
\end{thebibliography}
\end{document}
% END_OF_VERIFIED_REPORT::LSPV-5::PRIMECRT1
